\section{Error-bounded refinement along the width axis}
\label{sec:width}
Exact threshold integration leaves a second cost: evaluating every integer
width. We next bound the scalar average while omitting selected widths.
This is a different reduction from restricting thresholds. It uses the
geometry of expanded anomaly intervals and a probability-box (p-box)
expectation bound; the latter is established prior work
\cite{utkin2009,troffaes2011}.

\subsection{Where width monotonicity holds}
For consecutive original intervals ending at $e_i$ and starting at $s_{i+1}$,
their expansions first overlap at
\[
 c_i=2\left\lceil\frac{s_{i+1}-e_i}{2}\right\rceil .
\]
The strict overlap convention matters: merely adjacent expansions do not merge.
Partition $0,\ldots,L$ immediately before every distinct $c_i\le L$.
Within each resulting integer interval $[a,b]$, the number $K$ of expanded
segments is constant, each segment expands without merging, and each
threshold's hit count is nondecreasing.

\begin{lemma}[Coordinate order within a merge-free interval]
For a fixed threshold and $a\le w\le b$ in one such interval,
\[
 R_a\le R_w\le R_b,\quad
 \Prec_a\le\Prec_w\le\Prec_b,\quad
 F_a\ge F_w\ge F_b .
\]
At any fixed width, recall and FPR are nondecreasing as the threshold falls.
\end{lemma}
\begin{proof}
Let $O$ be the number of admitted original positives and $B_w$ the admitted
buffer mass. Both $O$ and the prediction count $C$ are fixed across widths,
while $B_w$ is nondecreasing. The unclipped overlap recall is
$(O+B_w)/(P_0+B_w/2)$, whose derivative in $B_w$ is
$(P_0-O/2)/(P_0+B_w/2)^2>0$. Multiplication by the nondecreasing
existence reward $H_w/K$ and clipping preserve order.
Precision is $(O+B_w)/C$. For FPR,
\[
 \frac{\partial F}{\partial B}
 =\frac{-(n-P_0)+(C-O)/2}{(n-P_0-B/2)^2}<0,
\]
since $C-O\le n-P_0$.
Across descending thresholds, original-positive admissions increase $O$;
other admissions increase $B$ and increase the false-positive numerator by
$1-a_i^{(w)}\ge0$ while decreasing its denominator. Thus both curve
coordinates are nondecreasing in threshold order. Ties enter together.
\end{proof}

\begin{theorem}[ROC width enclosure]
Within a merge-free interval, the exact ROC area is nondecreasing in width.
Consequently every omitted width has area in
$[\mathrm{ROC}_a,\mathrm{ROC}_b]$.
\end{theorem}
\begin{proof}
Use the common full threshold sequence, with the prescribed origin and
terminal point. Its trapezoidal area is coordinatewise nondecreasing in each
recall and nonincreasing in each interior FPR:
\[
 \frac{\partial A}{\partial R_j}=\frac{F_{j+1}-F_{j-1}}2\ge0,\qquad
 \frac{\partial A}{\partial F_j}=\frac{R_{j-1}-R_{j+1}}2\le0 .
\]
Linear interpolation between the endpoint coordinate sequences preserves
their threshold monotonicity. Apply the coordinate-order lemma along this
interpolation. Sparse event integration gives the same scalar area.
\end{proof}

Neither global ROC monotonicity nor PR-area monotonicity is assumed.
For example, with $n=20$, positive positions $\{0,2,19\}$ and scores
$(20,4,19,3,18,17,16,15,14,13,12,11,10,9,8,7,6,5,2,1)$,
ROC drops from $0.444444$ at width one to $0.405557$ at width two when
three expanded segments become two.
Even without a merge, scalar endpoint PR values need not enclose the
interior. Take $n=700$, positive positions $100,\ldots,199$ and $600$,
and initially zero scores. Assign score three to positions $100,\ldots,149$,
score one to $150,\ldots,199$ and $600$, and score two to
$0,\ldots,99$, $200,\ldots,549$, and $599$. The first merge occurs at
width 402. Direct evaluation gives PR $0.385206$ at width one,
$0.366313$ at width two, and $0.909134$ at width 400.
Thus the width-two area lies below both endpoint values on $[1,400]$.
The artifact verifies this witness independently.

\subsection{PR bounds without monotone precision}
Index the common descending support events by $j=1,\ldots,G$.
At every width, $q_j=R_j-R_{j-1}\ge0$ and $\sum_jq_j=1$, since the
last support event has recall one. Hence PR is an expectation
$\sum_j q_j\Prec_j$, with cumulative distribution $R_j$.
For $z\in(0,1)$, define the feasible quantile window
\[
 \ell(z)=\min\{j:R_{b,j}\ge z\},\qquad
 u(z)=\min\{j:R_{a,j}\ge z\}.
\]
The coordinate-order lemma implies that the quantile index for any
intermediate width belongs to $[\ell(z),u(z)]$.

\begin{theorem}[PR width enclosure]
Every width in $[a,b]$ has exact PR area between
\begin{align}
 L_{a,b}^{\mathrm{PR}}
 &=\int_0^1\min_{\ell(z)\le j\le u(z)}\Prec_{a,j}\,dz,\\
 U_{a,b}^{\mathrm{PR}}
 &=\int_0^1\max_{\ell(z)\le j\le u(z)}\Prec_{b,j}\,dz .
\end{align}
Both bounds can be computed in $O(G)$ time and space.
\end{theorem}
\begin{proof}
The quantile representation of $\sum_jq_j\Prec_j$ integrates the precision
at the quantile index over $z$. That index lies in the displayed window,
and its precision is between the corresponding endpoint precisions.
Taking the window minimum and maximum and integrating yields the bounds.
Breakpoints are the union of the two endpoint recall sequences. Merge them
in linear time. Both ends of the index window move only right; standard
monotone deques maintain the precision minimum and maximum in amortized
constant work per inserted or removed index. Each index is inserted and
removed at most once.
\end{proof}
These are extrema over a relaxed p-box model, not necessarily values
attained by an intermediate VUS width. This relaxation permits a valid bound
despite oscillating precision. We claim the VUS-specific coordinate order,
merge partition, and adaptive use of these bounds, not the general p-box
expectation construction.

\subsection{Adaptive evaluation and decisions}
Evaluate the endpoints of each merge-free interval. If an active interval
$[a,b]$ has $r=b-a-1$ unevaluated interior widths, add $rL_{a,b}$ and
$rU_{a,b}$ to the lower and upper sums, respectively. Add evaluated scalar
values exactly once, then divide by $L+1$. This encloses the target discrete
average; it does not replace it with a continuous-width integral.

Refine the interval with the largest contribution to uncertainty by
evaluating its midpoint and replacing it by two subintervals.
Stop once both metric intervals have half-width at most $\varepsilon$,
or when every width is evaluated. The returned estimate is the midpoint.
One model outranks another whenever its lower bound exceeds the other's
upper bound; overlapping intervals produce no ordering assertion.

With $m$ distinct evaluated widths, preparation, endpoint calculation,
linear p-box bounds and heap refinement take
$O(n\log(Q+1)+mQ+m\log m)$ time.
A fixed-size endpoint-curve cache retains at most eight curves; evicted
endpoints can be recomputed a constant number of times per refinement.
Space is $O(n+Q+m)$, with no width-by-support matrix.
The artifact records both distinct widths and actual curve evaluations.
The bound is output-sensitive to the requested precision and geometry;
neither $m\ll L$ nor a wall-clock improvement is guaranteed.

The theorems are real-arithmetic statements. The implementation uses
float64, clips ulp-scale recall drift, and adds an empirical $10^{-12}$
guard to adaptive endpoints. It is not a directed-rounding interval library.
Experiments independently audit enclosure within $10^{-10}$ and use a
$2\cdot10^{-10}$ separation margin for numerical ordering checks.
These allowances are reported separately from the requested approximation
half-width. Setting $\varepsilon=0$ evaluates every width.
